\documentclass[11pt,a4paper]{article}
\usepackage[margin=23mm]{geometry}
\usepackage{amsmath,amssymb,amsthm}
\setlength{\emergencystretch}{2em}
\newtheorem{theorem}{Theorem}
\newcommand{\gen}[1]{\langle#1\rangle}
\newcommand{\Alt}{\operatorname{Alt}}
\title{The first internal orbit at odd-degree equality}
\author{Human Compression project: L044 continuation}
\date{}
\begin{document}\maketitle
Let $G=\gen{\rho_0,\ldots,\rho_{r-1}}\le\Alt(\Omega)$
be transitive with a connected string C-marking and all
maximal omissions intransitive. Assume $n=2r+1$ and
the actual transitive six-point head on $D=\{1,\ldots,6\}$,
with the literal interface
\[
 \tau=\rho_3=(6\ b)\delta,\quad
 R=G_{\ge3}=\Alt(\Omega\setminus\{1,\ldots,5\}),\quad
 \sigma=\rho_4,\quad \sigma|_D=1.
\]
The involution $\delta$ fixes $D,b$. Put $H=G_{\ge5}$
and $K=G_{>5}$, and suppose $H$ fixes exactly $D\cup\{b\}$.
Use actual normalized assignments of two full ambient-eligible
edges per tail colour to a single $H$-orbit, as in the
acquired assignment theorem.

\begin{theorem}
Suppose the two nontrivial $H$-orbits are $O$ of size
$2q$ and $V$ of size $2s$, where $q\ge3$ is odd and
$s\ge1$. Assume $\delta$ preserves $O$ and acts freely
there, and that their assigned sets satisfy
\[
          5\in F_O,\qquad |F_O|=q-1,\qquad |F_V|=s-1.
\]
Then this configuration is impossible.
\end{theorem}
\begin{proof}
String commutation makes $\delta$ centralize $H$.
It therefore permutes actual $H$-orbits. Since it
preserves $O$ and fixes $D,b$, it also preserves $V$.
The actual pure-overlap argument gives
\[
 x=b^\sigma\ne b,\qquad \delta x=x,
 \qquad Kx=x,\qquad
 (6\ b\ x)^{\pm1}\in\gen{\tau,\sigma}.                 \tag{1}
\]
These are the acquired conclusions with natural $R$
and full ambient IP. They require no transitivity of
$G_{\ge4}$. The point $x$ is outside $D,b$ and cannot
belong to the free orbit $O$, so lies in $V$.
As $\delta|_V$ centralizes a transitive action and
fixes $x$, it fixes all of $V$.

The internal pair-forest equality, including the
first assigned colour, gives
\[
 F_O=\{5,\ldots,q+3\},\qquad
 I=F_V=\{q+4,\ldots,r-1\}.                             \tag{2}
\]
On the $q$ pairs of $\delta|_O$, the marks in $F_O$
are the successive transpositions of a path in this
order or its reverse. Every other tail mark has
trivial pair image. Thus an edge of a colour in $I$
on $O$ can only join a $\delta$-pair; it is already
joined by $\tau\in G_j$ and is ambient ineligible.
The seven fixed points carry no tail edges. All eligible
edges of the colours in $I$ are confined to $V$.

First suppose $s\ge2$, so $I$ is nonempty, and put
$T=G_I$. Since $K$ fixes $x$ but $H$ is transitive
on $V$, the mark $\rho_5$ must move $x$ to a distinct
point $y\in V$. Every mark in $I$ has index at least
$q+4\ge7$, so $T$ commutes with $\rho_5$. It fixes
$x$ as a subgroup of $K$ and hence also fixes $y$.
Moreover it acts nontrivially on $V$: each colour
in $I$ has two assigned eligible edges there.
Consequently $o_T(V)\ge3$, counting the two fixed
points and at least one nontrivial orbit.
The complete acquired localization theorem gives
\[
             2s-2=2|I|\le |V|-o_T(V)\le2s-3,
\]
a contradiction. This argument neither identifies a
selected equality forest nor discards its possible
square components.

It remains that $s=1$, so $V=\{x,y\}$ and $I$ is
empty. The group $K$ fixes both points of $V$.
Removing colour $5$ from the pair path in (2) leaves
one isolated pair and the natural $S_{q-1}$ path on
the remaining pairs. In particular any $K$-fixed
point of $O$ must belong to that isolated pair;
$q-1\ge2$.

The groups $H$ and $\gen\tau$ preserve $O$. Since
$R=\gen{\tau,\sigma,H}$ is transitive on its natural
support, $\sigma$ cannot preserve $O$. It fixes $D$
and swaps $b,x$. The only remaining point outside
$O$ is $y$, so it must swap $y$ with some $z\in O$.
Commutation with $K$ makes $z$ fixed by $K$.
The mate $z'=z^\delta$ is also fixed by $K$, and
these are exactly the points of the isolated pair.
Thus the full $K$-fixed set is
\[
                       D\cup\{b,x,y,z,z'\}.
\]
The involution $\sigma$ preserves this set, fixes
$D$, and has already paired $b,x$ and $y,z$.
It must fix $z'$.

The product $\tau\sigma$ is now a three-cycle on
each of the two disjoint invariant sets
$\{6,b,x\}$ and $\{y,z,z'\}$. Every odd-order
element of the dihedral group $\gen{\tau,\sigma}$
belongs to its cyclic rotation subgroup. A power
trivial on the second triple is divisible by three
and is trivial on the first. This excludes the pure
three-cycle asserted in (1), the final contradiction.
\end{proof}

\paragraph{Attribution and boundary.}
The normalized budget, internal pair-path equality,
pure head overlap and complete ambient localization
are acquired inputs. The new step is the second
$T$-fixed anchor supplied by $\rho_5$, followed by
the literal two-point residual. The proof retains
all square components and uses no projected IP or
signed lift classification. Other seven-fixed-point
budgets, pairs split between different $H$-orbits,
global separator transfer and the complete L044
recovery remain separate. No new alternating rank
induction, computation, formal declaration, original
or question credit is claimed.
\end{document}
